\chapter{晶体点群特征标：生成规则与常用表}
\label{app:crystal-character-tables}

点群的几何分类建立以后，下一个问题是如何把对称操作转换成可计算的表示。特征标表同时编码群的共轭类、不可约表示，以及常见物理基函数在这些表示中的归属。本章从\chapref{chap:point-space-groups}的晶体点群出发，先建立生成规则，再整理常用显式表；这些表示数据可直接用于\chapref{chap:group-theory-quantum}的量子力学计算。这样三十二种晶体点群的表就不再是三十二份彼此孤立的记忆材料。即使忘记某个数字，也能从群结构重新算出来。

第一次查特征标表时，建议先按“列、行、表内数字、基函数”四层读取。表的每一\emph{列}通常对应一个共轭类，列标题前的系数如 $2C_3$ 表示这一类含两个操作；每一\emph{行}对应一个不可约表示，行名 $A,B,E,T$ 在常用 Mulliken 记号中通常分别表示一维、一维、二维、三维类型。表格中央的数字是特征标，也就是该类任一群元表示矩阵的迹，因此恒等元 $E$ 那一列的特征标必等于表示维数。表边上的 $x,y,z$、$R_x,R_y,R_z$、$x^2-y^2$ 等则是在告诉你哪些常见函数或物理量按这一行变换。查表时先看恒等元列判断维数，再看目标物理量所在行，通常比从左到右死读整张表有效得多。

需要首先澄清一个常见约定。全书默认在复数域上讨论不可约表示，因此例如循环群 $C_3$ 有三个一维复不可约表示。晶体学表格为了使用实坐标基，且常默认体系具有时间反演对称性，会把一对互为复共轭的一维表示组合成一个二维实基组，并沿用字母 $E$ 标记。这个“二维”是实表示或物理配对后的维数，并不意味着原来的两个复一维表示在纯点群意义下变成了同一个不可约表示。以下公式先坚持复表示，再在需要时说明实基组合。

\section{三十二种晶体点群的索引}

晶体制约定理只允许 $1,2,3,4,6$ 重旋转轴。将\chapref{chap:point-space-groups}得到的有限点群与这一限制相交，恰好得到三十二种晶体点群：

\begin{table}[htbp]
\centering
\caption{三十二种晶体点群按晶系分类。}
\label{tab:32-crystal-point-groups-app}
\begin{tabular}{ccl}
\toprule
晶系 & 数目 & Schoenflies 符号 \\
\midrule
三斜 & 2 & $C_1,\ C_i$ \\
单斜 & 3 & $C_2,\ C_s,\ C_{2h}$ \\
正交 & 3 & $D_2,\ C_{2v},\ D_{2h}$ \\
四方 & 7 & $C_4,\ S_4,\ C_{4h},\ D_4,\ C_{4v},\ D_{2d},\ D_{4h}$ \\
三方 & 5 & $C_3,\ S_6,\ D_3,\ C_{3v},\ D_{3d}$ \\
六方 & 7 & $C_6,\ C_{3h},\ C_{6h},\ D_6,\ C_{6v},\ D_{3h},\ D_{6h}$ \\
立方 & 5 & $T,\ T_h,\ O,\ T_d,\ O_h$ \\
\bottomrule
\end{tabular}
\end{table}

这个表不是三十二个彼此独立的对象。大量点群可由循环群、二面体群、反演群和镜面对称做直积或半直积构造，因此特征标也可以系统生成。

\section{循环族与二面体族}

\subsection{\texorpdfstring{循环群 $C_n$}{循环群 Cn}}

令 $C_n=\langle r\mid r^n=e\rangle$，并令
\[
\omega_n=\ee^{2\pi\ii/n}.
\]
\begin{proposition}[循环群的全部复不可约特征标]
$C_n$ 的共轭类均为单元素类。其全部复不可约表示为以下 $n$ 个一维表示：
\begin{equation}
\rho_m(r^k)=\omega_n^{mk},
\qquad m=0,1,\dots,n-1.
\label{eq:cn-characters-app}
\end{equation}
因此
\[
\chi_m(r^k)=\omega_n^{mk}.
\]
\end{proposition}
\begin{proof}
$C_n$ 为 Abel 群，Schur 引理使每个不可约表示中的 $\rho(r)$ 都为标量；不可约性进一步要求表示空间为一维。关系 $r^n=e$ 给出 $\rho(r)^n=1$，故 $\rho(r)=\omega_n^m$。不同 $m$ 的正交关系化为离散 Fourier 正交：
\[
\frac1n\sum_{k=0}^{n-1}\omega_n^{(m'-m)k}=\delta_{mm'}.
\]
这是有限等比数列求和的结果，且 $n$ 个表示的维数平方和为 $n$，故没有遗漏。
\end{proof}
所以 $C_n$ 的特征标表本质上就是 $n$ 点离散 Fourier 变换矩阵。

\begin{proposition}[共轭特征标的实基组合]
对 $m\not\equiv-m\pmod n$，$m$ 与 $n-m$ 两个复共轭表示合成一个二维实不可约表示，其特征标为
\begin{equation}
\chi^{\RR}_m(r^k)=2\cos\frac{2\pi mk}{n}.
\label{eq:cn-real-character}
\end{equation}
若 $m\equiv-m\pmod n$，相应实表示已经是一维，不能再按二维不可约表示计数。
\end{proposition}
\begin{proof}
在二维实空间取基
\[
\cos(m\phi),\qquad \sin(m\phi),
\]
则 $r^k$ 的实矩阵为转角 $2\pi mk/n$ 的二维转动，迹等于两个复特征标之和。若 $m\not\equiv-m$，转动矩阵没有实不变直线，故实表示不可约；例外恰为 $m=0$，以及偶数 $n$ 时的 $m=n/2$。
\end{proof}
这正是很多晶体学表中 $E$ 型二维行的来源。

\subsection{\texorpdfstring{二面体群 $D_n$ 与 $C_{nv}$}{二面体群 Dn 与 Cnv}}

采用群表示
\[
D_n=\langle r,s\mid r^n=e,\ s^2=e,\ srs=r^{-1}\rangle,
\qquad |D_n|=2n.
\]
几何上，$r$ 是主轴转动，$s$ 是一条与主轴垂直的二重轴；若把 $s$ 解释成包含主轴的镜面，则同一个抽象群就是 $C_{nv}$。因此 $D_n$ 与 $C_{nv}$ 的抽象特征标结构相同，区别只在共轭类的物理标签与基函数解释。

\begin{theorem}[二面体群的复不可约表示]
若 $n$ 为奇数，除两个一维表示外，有 $(n-1)/2$ 个二维表示 $E_m$：
\begin{equation}
\chi_{E_m}(r^k)=2\cos\frac{2\pi mk}{n},
\qquad
\chi_{E_m}(sr^k)=0,
\qquad
m=1,\dots,\frac{n-1}{2}.
\label{eq:dn-odd-char}
\end{equation}
两个一维表示由 $r\mapsto1$、$s\mapsto\pm1$ 给出。

若 $n$ 为偶数，$r$ 还可在一维表示中取 $-1$，所以共有四个一维表示；二维表示为 $m=1,\dots,n/2-1$，其特征标仍由 \eqnrefs{eq:dn-odd-char} 给出。此时两类“反射型”元素 $sr^{2k}$ 与 $sr^{2k+1}$ 不再共轭，这正对应 $D_n$ 中两类横向二重轴，或 $C_{nv}$ 中 $\sigma_v$ 与 $\sigma_d$ 的区分。
\end{theorem}
\begin{proof}
在一维表示中，$srs=r^{-1}$ 要求 $\rho(r)^2=1$；再结合 $\rho(r)^n=1$、$\rho(s)^2=1$，得到奇数 $n$ 的两组和偶数 $n$ 的四组符号选择。

二维表示可显式取为
\[
\rho_m(r)=\begin{pmatrix}\omega_n^m&0\\0&\omega_n^{-m}\end{pmatrix},
\qquad
\rho_m(s)=\begin{pmatrix}0&1\\1&0\end{pmatrix}.
\]
这些矩阵满足生成关系；$s$ 交换 $r$ 的两条不同本征线，故不存在一维不变子空间。矩阵的迹分别给出 \eqnrefs{eq:dn-odd-char} 中的余弦与零；$m$ 与 $n-m$ 等价，所列范围恰好去除重复。维数平方和为
\[
2+4\frac{n-1}{2}=2n\quad(n\text{ 为奇数}),\qquad
4+4\left(\frac n2-1\right)=2n\quad(n\text{ 为偶数}),
\]
因而表示列表完备。最后，$r^\ell(sr^k)r^{-\ell}=sr^{k-2\ell}$；奇数 $n$ 时所有 $k$ 相通，偶数 $n$ 时奇偶性不变，得到两类反射型元素。
\end{proof}

\begin{example}[$C_{3v}$]
由奇数阶二面体群的一般特征标公式构造 $C_{3v}$ 的特征标表，并列出常见坐标基。

\textbf{解。} $C_{3v}\cong D_3$，三个共轭类为 $E,2C_3,3\sigma_v$。应用上面的奇数 $n$ 公式，得到
\begin{center}
\begin{tabular}{c|ccc|l}
\toprule
 & $E$ & $2C_3$ & $3\sigma_v$ & 常见基 \\
\midrule
$A_1$ & $1$ & $1$ & $1$ & $z,\ x^2+y^2,\ z^2$ \\
$A_2$ & $1$ & $1$ & $-1$ & $R_z$ \\
$E$   & $2$ & $-1$ & $0$ & $(x,y)$，$(R_x,R_y)$，$(x^2-y^2,2xy)$ \\
\bottomrule
\end{tabular}
\end{center}
这里的二维 $E$ 正是公式 $2\cos(2\pi/3)=-1$ 的结果，而不是需要单独记忆的数字。
\end{example}

\begin{example}[$D_4$]
由偶数阶二面体群的一般结构构造 $D_4$ 的特征标表，并核对各共轭类的大小。

\textbf{解。} $D_4$ 的类可排列为 $E,2C_4,C_2,2C_2',2C_2''$。由偶数 $n$ 的一般结构得到
\begin{center}
\begin{tabular}{c|rrrrr}
\toprule
 & $E$ & $2C_4$ & $C_2$ & $2C_2'$ & $2C_2''$ \\
\midrule
$A_1$ & 1& 1& 1& 1& 1\\
$A_2$ & 1& 1& 1&-1&-1\\
$B_1$ & 1&-1& 1& 1&-1\\
$B_2$ & 1&-1& 1&-1& 1\\
$E$   & 2& 0&-2& 0& 0\\
\bottomrule
\end{tabular}
\end{center}
这张表随后可通过加入反演或镜面宇称生成 $D_{4h}$、$D_{2d}$ 等常用四方群的大部分结构。
\end{example}

\section{\texorpdfstring{反演与 $g/u$ 宇称}{反演与 g/u 宇称}}

\begin{proposition}[中心反演的宇称加倍]
若一个点群可写成
\[
G=K\times C_i,
\qquad C_i=\{E,I\},
\]
其中反演 $I$ 与 $K$ 的全部元素交换，则不可约表示就是 $K$ 的不可约表示与 $C_i$ 的两个一维表示做张量积。后者的特征标为
\[
\chi_g(E)=1,\quad \chi_g(I)=1;
\qquad
\chi_u(E)=1,\quad \chi_u(I)=-1.
\]
因此每个 $K$ 的不可约表示 $\Gamma$ 自动产生
\[
\Gamma_g=\Gamma\otimes g,
\qquad
\Gamma_u=\Gamma\otimes u.
\]
在 $K$ 部分两者特征标相同，在反演陪集 $IK$ 上只差整体负号。$D_{2h},D_{4h},D_{6h},T_h,O_h$ 的特征标表都可以用这一原则大幅简化，而不是从零开始求解。
\end{proposition}
\begin{proof}
在不可约表示中，中心元 $I$ 由 Schur 引理表示为标量，又因 $I^2=E$，只能取 $+1$ 或 $-1$。固定这一符号后，不变子空间完全由限制到 $K$ 的表示决定，故该限制不可约；反过来，任一 $K$ 的不可约表示乘上此符号便给出 $G$ 的不可约表示。其特征标满足
\[
\chi_{\Gamma_g}(Ik)=\chi_\Gamma(k),\qquad
\chi_{\Gamma_u}(Ik)=-\chi_\Gamma(k),\qquad k\in K,
\]
这也直接说明两份表的符号关系。
\end{proof}

对物理基函数，极矢量 $(x,y,z)$ 在反演下变号，故为 $u$；轴矢量 $(R_x,R_y,R_z)$ 在反演下不变，故为 $g$。二次函数 $x^2,y^2,z^2,xy,yz,zx$ 都是偶宇称。只要再知道它们在纯转动子群下怎样变换，就能立即给出完整 $g/u$ 标签。

\section{\texorpdfstring{立方族：$T,O,T_d,O_h$}{立方族：T, O, Td, Oh}}

\subsection{\texorpdfstring{正八面体旋转群 $O$}{正八面体旋转群 O}}

$O$ 是立方体或正八面体的纯转动群，阶为 $24$。按
\[
E,\qquad 8C_3,\qquad 3C_2,\qquad 6C_4,\qquad 6C_2'
\]
排列共轭类，其五个不可约表示为
\begin{table}[htbp]
\centering
\caption{$O$ 群特征标表。}
\label{tab:O-character}
\begin{tabular}{c|rrrrr|l}
\toprule
 & $E$ & $8C_3$ & $3C_2$ & $6C_4$ & $6C_2'$ & 常见基 \\
\midrule
$A_1$ & 1& 1& 1& 1& 1 & $x^2+y^2+z^2$\\
$A_2$ & 1& 1& 1&-1&-1 & \\
$E$   & 2&-1& 2& 0& 0 & $(2z^2-x^2-y^2,x^2-y^2)$\\
$T_1$ & 3& 0&-1& 1&-1 & $(x,y,z)$ 或 $(R_x,R_y,R_z)$，视极/轴矢量而定\\
$T_2$ & 3& 0&-1&-1& 1 & $(xy,yz,zx)$\\
\bottomrule
\end{tabular}
\end{table}

维数平方和
\[
1^2+1^2+2^2+3^2+3^2=24
\]
首先检验了表的完备性。把 $O$ 与反演群直积即可得到 $O_h$ 的十个不可约表示
\[
A_{1g},A_{2g},E_g,T_{1g},T_{2g},
A_{1u},A_{2u},E_u,T_{1u},T_{2u}.
\]
\chapref{chap:group-theory-quantum}中立方晶场下 $d$ 轨道的分裂
\[
D^{(l=2)}\downarrow O_h=E_g\oplus T_{2g}
\]
正是由 \tabref{tab:O-character} 中二次基函数一列读出的。

\subsection{\texorpdfstring{正四面体群 $T$ 与 $T_d$}{正四面体群 T 与 Td}}

$T\cong A_4$，阶为 $12$。在复表示理论中，它有三个一维表示与一个三维表示。将类排列为 $E,4C_3,4C_3^2,3C_2$，记 $\omega=\ee^{2\pi\ii/3}$，则
\begin{center}
\begin{tabular}{c|rrrr}
\toprule
 & $E$ & $4C_3$ & $4C_3^2$ & $3C_2$\\
\midrule
$A$   &1&1&1&1\\
$A'$  &1&$\omega$&$\omega^2$&1\\
$A''$ &1&$\omega^2$&$\omega$&1\\
$T$   &3&0&0&-1\\
\bottomrule
\end{tabular}
\end{center}
$A'$ 与 $A''$ 互为复共轭。若只使用实基或再加入时间反演，它们经常作为一个二维 $E$ 型基组共同出现；这正是本章开头所强调的“复不可约表示”与“晶体学实基约定”的区别。

$T_d$ 的类可取 $E,8C_3,3C_2,6S_4,6\sigma_d$，常用表为
\begin{center}
\begin{tabular}{c|rrrrr|l}
\toprule
 & $E$&$8C_3$&$3C_2$&$6S_4$&$6\sigma_d$&常见基\\
\midrule
$A_1$&1&1&1&1&1&$x^2+y^2+z^2$\\
$A_2$&1&1&1&-1&-1&\\
$E$&2&-1&2&0&0&$(2z^2-x^2-y^2,x^2-y^2)$\\
$T_1$&3&0&-1&1&-1&$(R_x,R_y,R_z)$\\
$T_2$&3&0&-1&-1&1&$(x,y,z)$，$(xy,yz,zx)$\\
\bottomrule
\end{tabular}
\end{center}

\section{由基函数重建特征标}

\begin{proposition}[不变函数空间的特征标重建]
若线性无关函数 $f_1,\dots,f_d$ 的张成空间在群作用 $U$ 下闭合，定义
\[
U(g)f_j=\sum_i f_i D_{ij}(g),
\]
则
\[
\chi(g)=\Tr D(g).
\]
若换一组基函数，所得特征标不变。
\end{proposition}
\begin{proof}
线性无关性使系数矩阵唯一。由 $U(g)U(h)=U(gh)$ 得到 $D(g)D(h)=D(gh)$，所以 $D$ 是表示。换基 $f'=fS$ 给出 $D'(g)=S^{-1}D(g)S$；迹在相似变换下不变。
\end{proof}

\begin{example}[$C_{4v}$ 中的平面坐标表示]
让 $C_{4v}$ 作用在平面坐标 $(x,y)$ 上，计算各类特征标并识别所得二维表示。

\textbf{解。} 在 $C_{4v}$ 中，基 $(x,y)$ 在 $C_4$ 下变换为
\[
\begin{pmatrix}x\\y\end{pmatrix}
\mapsto
\begin{pmatrix}0&-1\\1&0\end{pmatrix}
\begin{pmatrix}x\\y\end{pmatrix},
\]
故 $\chi(C_4)=0$；在 $C_2$ 下矩阵为 $-I_2$，故 $\chi(C_2)=-2$；在任一竖直镜面下本征值为 $+1,-1$，故迹为 $0$。于是无需查表就恢复了 $C_{4v}$ 的二维 $E$ 行
\[
(2,0,-2,0,0).
\]
\end{example}

这个方法把前面的表格组织成同一种计算程序。真正进入科研时，表格可以由数据库提供，但“为什么某组轨道、位移、振动模式或算符属于某个不可约表示”仍要靠这样的表示构造来判断。
